\documentclass[12pt,a4paper]{article}
\usepackage[margin=2.5cm]{geometry}
\usepackage{float}
\usepackage{tabularx}
\usepackage{booktabs}
\usepackage{colortbl}
\usepackage{amsmath}
\usepackage{amssymb}
\usepackage{enumitem}
\usepackage{hyperref}
\usepackage{fancyhdr}
\setlength{\headheight}{15pt}
\usepackage[dvipsnames,table]{xcolor}
\usepackage{tcolorbox}
\tcbuselibrary{skins, breakable}
\usepackage{titlesec}

% Section styling
\titleformat{\section}{\color{Maroon}\Large\bfseries}{\thesection}{1em}{}[\color{Maroon}\titlerule]
\titleformat{\subsection}{\color{BrickRed}\large\bfseries}{\thesubsection}{1em}{}

% Custom boxes
\newtcolorbox{sqlbox}[1][]{
  colback=blue!2,
  colframe=Maroon,
  fonttitle=\bfseries,
  title=#1,
  boxrule=1pt,
  arc=3pt,
  breakable
}

\newtcolorbox{metabox}[1][]{
  colback=green!2,
  colframe=OliveGreen,
  fonttitle=\bfseries,
  title=#1,
  boxrule=0.5pt,
  arc=2pt,
  breakable
}

\pagestyle{fancy}
\fancyhf{}
\rhead{ISE 305 -- HW4}
\lhead{Aydogan \& D\"{o}nmez}
\cfoot{\thepage}

\title{\textbf{ISE 305 -- Database Systems}\\Homework 4: Complex SQL Queries\\[6pt]\large Spring 2026 -- Instructor: Ali \c{C}akmak}
\author{Y\d{i}\u{g}it Aydogan \quad 150230208\\Emir Mete D\"{o}nmez \quad 150220205}
\date{May 18, 2026}

\begin{document}
\maketitle

% ══════════════════════════════════════════════════════════════════════════════
\section{Part I: Database Schema}
% ══════════════════════════════════════════════════════════════════════════════

The NutriQuery database consists of \textbf{8 tables} storing USDA FoodData Central and
OpenFoodFacts data, brand metadata, nutritional profiles, health/allergen attributes, and
machine-learning prediction results.

\subsection{Schema Listing}

\begin{table}[H]\centering\small
\resizebox{\textwidth}{!}{%
\begin{tabular}{llllp{4.5cm}}
\toprule
\textbf{Table} & \textbf{Column} & \textbf{Data Type} & \textbf{Constraint} & \textbf{Description} \\
\midrule
\rowcolor{gray!10}
\multicolumn{5}{l}{\textbf{Brands}} \\
& brand\_id    & INT          & \textbf{PK} (IDENTITY) & Surrogate primary key \\
& brand\_name  & VARCHAR(255) & NOT NULL               & Brand display name \\
& brand\_owner & VARCHAR(255) &                        & Parent company or owner name \\
\midrule
\rowcolor{gray!10}
\multicolumn{5}{l}{\textbf{FOOD\_CATEGORY}} \\
& category\_id   & INT          & \textbf{PK} (IDENTITY) & Surrogate primary key \\
& category\_name & VARCHAR(255) & NOT NULL, UNIQUE       & Food group name (e.g., Snacks, Beverages) \\
\midrule
\rowcolor{gray!10}
\multicolumn{5}{l}{\textbf{DATA\_TYPE}} \\
& type\_id   & INT          & \textbf{PK} (IDENTITY) & Surrogate primary key \\
& type\_name & VARCHAR(100) & NOT NULL, UNIQUE       & USDA data source (e.g., SR Legacy, Branded) \\
\midrule
\rowcolor{gray!10}
\multicolumn{5}{l}{\textbf{Foods}} \\
& fdc\_id      & INT          & \textbf{PK}            & USDA FoodData Central identifier \\
& brand\_id    & INT          & \textbf{FK} $\to$ Brands(brand\_id) & Referenced brand (nullable) \\
& category\_id & INT          & \textbf{FK} $\to$ FOOD\_CATEGORY(category\_id) & Food category (nullable) \\
& type\_id     & INT          & \textbf{FK} $\to$ DATA\_TYPE(type\_id) & Data source type (nullable) \\
& food\_name   & VARCHAR(500) & NOT NULL               & Descriptive food name \\
\midrule
\rowcolor{gray!10}
\multicolumn{5}{l}{\textbf{Nutrition\_Metrics}} \\
& nutrition\_id & INT  & \textbf{PK} (IDENTITY) & Surrogate primary key \\
& fdc\_id       & INT  & NOT NULL, \textbf{FK} $\to$ Foods(fdc\_id) & One-to-one with Foods \\
& calories     & FLOAT &                       & Energy (kcal per 100g) \\
& protein\_g   & FLOAT &                       & Protein content (g per 100g) \\
& fat\_g       & FLOAT &                       & Total fat (g per 100g) \\
& carbs\_g     & FLOAT &                       & Carbohydrates (g per 100g) \\
& sodium\_mg   & FLOAT &                       & Sodium (mg per 100g) \\
\midrule
\rowcolor{gray!10}
\multicolumn{5}{l}{\textbf{HEALTH\_SCORE}} \\
& score\_id        & INT         & \textbf{PK} (IDENTITY) & Surrogate primary key \\
& fdc\_id          & INT         & NOT NULL, UNIQUE, \textbf{FK} $\to$ Foods(fdc\_id) & One-to-one with Foods \\
& health\_score    & FLOAT       &                        & Composite health score (0--100) \\
& nutriscore\_grade & VARCHAR(50) &                       & Nutri-Score grade (A--E) \\
& nova\_group      & INT         &                        & NOVA processing group (1--4) \\
\midrule
\rowcolor{gray!10}
\multicolumn{5}{l}{\textbf{ALLERGEN\_PROFILE}} \\
& allergen\_id    & INT  & \textbf{PK} (IDENTITY)  & Surrogate primary key \\
& fdc\_id         & INT  & NOT NULL, UNIQUE, \textbf{FK} $\to$ Foods(fdc\_id) & One-to-one with Foods \\
& contains\_gluten & BIT & DEFAULT 0               & Gluten presence flag \\
& contains\_dairy  & BIT & DEFAULT 0               & Dairy presence flag \\
\midrule
\rowcolor{gray!10}
\multicolumn{5}{l}{\textbf{ML\_Predictions}} \\
& prediction\_id       & INT       & \textbf{PK} (IDENTITY) & Surrogate primary key \\
& fdc\_id              & INT       & NOT NULL, \textbf{FK} $\to$ Foods(fdc\_id) & Predicted food item \\
& predicted\_nutriscore & VARCHAR(50) &                     & ML-predicted Nutri-Score grade \\
& predicted\_nova      & INT       &                        & ML-predicted NOVA group \\
& confidence\_score    & FLOAT     &                        & Softmax confidence (0--1) \\
& prediction\_date     & DATETIME  & DEFAULT GETDATE()      & Timestamp of prediction \\
\bottomrule
\end{tabular}%
}
\caption{NutriQuery database schema --- 8 tables. Primary keys are in \textbf{bold}; foreign keys are annotated with their target relation. UNIQUE constraints on \texttt{fdc\_id} in Nutrition\_Metrics, HEALTH\_SCORE, and ALLERGEN\_PROFILE enforce strict 1:1 relationships.}
\label{tab:schema}
\end{table}

\subsection{Schema Design Summary}
\begin{itemize}
  \item \textbf{Lookup Tables:} \texttt{FOOD\_CATEGORY} and \texttt{DATA\_TYPE} store category and data-source names independently of \texttt{Foods}, avoiding duplication. Category and type names are resolved through foreign keys \texttt{category\_id} and \texttt{type\_id}.
  \item \textbf{1:1 Enforcement:} \texttt{Nutrition\_Metrics}, \texttt{HEALTH\_SCORE}, and \texttt{ALLERGEN\_PROFILE} each hold a UNIQUE foreign key on \texttt{fdc\_id}, guaranteeing that every food has at most one profile in each table.
  \item \textbf{1:N Relationships:} \texttt{ML\_Predictions} is a 1:N child of \texttt{Foods}, allowing multiple predictions per food across retraining runs. \texttt{Brands}, \texttt{FOOD\_CATEGORY}, and \texttt{DATA\_TYPE} are each 1:N parents of \texttt{Foods}.
\end{itemize}


% ══════════════════════════════════════════════════════════════════════════════
\section{Part II: Complex SQL Queries}
% ══════════════════════════════════════════════════════════════════════════════

We present \textbf{6 complex SQL queries} that drive the core analytical features of the NutriQuery
web application. Each query joins at least 4 tables and satisfies at least one of the required
complexity criteria (JOIN $\geq$3 tables, OUTER JOIN, GROUP BY, or nested SUBQUERY).
All queries were validated against the live MSSQL database.

% ── Query 1 ────────────────────────────────────────────────────────────────────
\subsection{Query 1: Full Food Profile Retrieval}

\begin{sqlbox}[SQL Query]
\small
\begin{verbatim}
SELECT
    f.fdc_id,
    f.food_name,
    fc.category_name AS food_category,
    dt.type_name     AS data_type,
    b.brand_name,
    b.brand_owner,
    n.calories,
    n.protein_g,
    n.fat_g,
    n.carbs_g,
    n.sodium_mg,
    hs.health_score,
    hs.nutriscore_grade,
    hs.nova_group,
    ap.contains_gluten,
    ap.contains_dairy
FROM Foods f
LEFT JOIN Brands b             ON f.brand_id    = b.brand_id
LEFT JOIN FOOD_CATEGORY fc     ON f.category_id = fc.category_id
LEFT JOIN DATA_TYPE dt         ON f.type_id     = dt.type_id
LEFT JOIN Nutrition_Metrics n  ON f.fdc_id      = n.fdc_id
LEFT JOIN HEALTH_SCORE hs      ON f.fdc_id      = hs.fdc_id
LEFT JOIN ALLERGEN_PROFILE ap  ON f.fdc_id      = ap.fdc_id
WHERE f.fdc_id = @fdc_id;
\end{verbatim}
\end{sqlbox}

\begin{metabox}[Query Translation]
\begin{itemize}
  \item \textbf{English Meaning:} Retrieves the complete profile of a single food item identified
    by its USDA \texttt{fdc\_id}. It assembles brand metadata (via Brands), category and data source
    labels (via FOOD\_CATEGORY and DATA\_TYPE lookup tables), full nutrition facts, health scores,
    and allergen flags into one unified result row. LEFT JOINs ensure that foods with missing
    brand, category, or health data are still returned rather than silently dropped.
  \item \textbf{Project Feature:} Powers the \textbf{``Lookup by FDC ID''} search bar in the Product
    Registry panel on the NutriQuery home page. The frontend calls \texttt{GET /foods/\{fdc\_id\}}
    and renders the full nutritional breakdown.
  \item \textbf{Criteria Met:}
    \begin{enumerate}[label=\textbf{\arabic*}]
      \item JOIN $\geq$3 tables --- joins \textbf{Foods}, \textbf{Brands}, \textbf{FOOD\_CATEGORY},
        \textbf{DATA\_TYPE}, \textbf{Nutrition\_Metrics}, \textbf{HEALTH\_SCORE}, and
        \textbf{ALLERGEN\_PROFILE} (7 tables).
      \item OUTER JOIN --- uses six LEFT JOINs to preserve food records even when related
        rows in lookup or profile tables are absent.
    \end{enumerate}
\end{itemize}
\end{metabox}

\newpage

% ── Query 2 ────────────────────────────────────────────────────────────────────
\subsection{Query 2: Multi-Constraint Range Query}

\begin{sqlbox}[SQL Query]
\small
\begin{verbatim}
SELECT
    f.fdc_id,
    f.food_name,
    fc.category_name AS food_category,
    b.brand_name,
    n.calories,
    n.protein_g,
    n.fat_g,
    n.carbs_g,
    n.sodium_mg,
    hs.health_score,
    hs.nutriscore_grade
FROM Foods f
JOIN FOOD_CATEGORY fc        ON f.category_id = fc.category_id
JOIN Nutrition_Metrics n     ON f.fdc_id      = n.fdc_id
JOIN HEALTH_SCORE hs         ON f.fdc_id      = hs.fdc_id
LEFT JOIN Brands b           ON f.brand_id    = b.brand_id
WHERE hs.health_score >= @min_health_score
  AND n.sodium_mg    <= @max_sodium
  AND n.carbs_g      <= @max_carbs
ORDER BY hs.health_score DESC;
\end{verbatim}
\end{sqlbox}

\begin{metabox}[Query Translation]
\begin{itemize}
  \item \textbf{English Meaning:} Finds foods that simultaneously satisfy three user-defined
    nutritional thresholds: a minimum health score, a maximum sodium level, and a maximum
    carbohydrate level. Results are ranked from healthiest to least healthy. INNER JOINs are
    used for nutrition, health, and category data because a food without these metrics is
    irrelevant to a range-filtering use case.
  \item \textbf{Project Feature:} Drives the \textbf{Range Query} panel in the Nutrition
    Analytics section. Users adjust sliders for health score, sodium, and carbs; the frontend
    calls \texttt{GET /queries/range} and renders the filtered list in an AG Grid table.
  \item \textbf{Criteria Met:}
    \begin{enumerate}[label=\textbf{\arabic*}]
      \item JOIN $\geq$3 tables --- joins \textbf{Foods}, \textbf{FOOD\_CATEGORY},
        \textbf{Nutrition\_Metrics}, \textbf{HEALTH\_SCORE}, and \textbf{Brands}
        (5 tables).
      \item OUTER JOIN --- LEFT JOIN on Brands allows foods without a brand assignment
        to still appear in results.
    \end{enumerate}
\end{itemize}
\end{metabox}

% ── Query 3 ────────────────────────────────────────────────────────────────────
\subsection{Query 3: Category Nutritional Aggregation}

\begin{sqlbox}[SQL Query]
\small
\begin{verbatim}
SELECT
    fc.category_name AS food_category,
    AVG(n.calories)  AS avg_calories,
    AVG(n.protein_g) AS avg_protein,
    AVG(n.fat_g)     AS avg_fat,
    AVG(n.carbs_g)   AS avg_carbs,
    COUNT(f.fdc_id)  AS item_count
FROM Foods f
JOIN FOOD_CATEGORY fc     ON f.category_id = fc.category_id
JOIN Nutrition_Metrics n  ON f.fdc_id      = n.fdc_id
WHERE fc.category_name = @category
GROUP BY fc.category_name;
\end{verbatim}
\end{sqlbox}

\begin{metabox}[Query Translation]
\begin{itemize}
  \item \textbf{English Meaning:} Computes aggregate nutritional statistics --- average calories,
    protein, fat, and carbohydrates --- for all food items belonging to a user-selected category.
    The COUNT function reports the total number of items in that category with complete nutrition
    data. Renaming a category requires only a single-row update in FOOD\_CATEGORY
    with no cascading changes to Foods.
  \item \textbf{Project Feature:} Populates the \textbf{Category Aggregation} dashboard in the
    Nutrition Analytics section. The user picks a category from a dropdown (populated by
    \texttt{GET /categories/}), and the frontend renders five summary stat cards.
  \item \textbf{Criteria Met:}
    \begin{enumerate}[label=\textbf{\arabic*}]
      \item JOIN $\geq$3 tables --- joins \textbf{Foods}, \textbf{FOOD\_CATEGORY}, and
        \textbf{Nutrition\_Metrics} (3 tables).
      \setcounter{enumi}{2}
      \item GROUP BY --- uses \texttt{GROUP BY fc.category\_name} with four \texttt{AVG()}
        aggregates and \texttt{COUNT()}.
    \end{enumerate}
\end{itemize}
\end{metabox}

\newpage

% ── Query 4 ────────────────────────────────────────────────────────────────────
\subsection{Query 4: Dietary Restriction Filtering}

\begin{sqlbox}[SQL Query]
\small
\begin{verbatim}
SELECT
    f.fdc_id,
    f.food_name,
    fc.category_name AS food_category,
    b.brand_name,
    n.calories,
    n.protein_g,
    n.fat_g,
    n.carbs_g,
    hs.health_score,
    hs.nutriscore_grade
FROM Foods f
JOIN FOOD_CATEGORY fc       ON f.category_id = fc.category_id
JOIN ALLERGEN_PROFILE ap    ON f.fdc_id      = ap.fdc_id
LEFT JOIN HEALTH_SCORE hs   ON f.fdc_id      = hs.fdc_id
LEFT JOIN Nutrition_Metrics n ON f.fdc_id    = n.fdc_id
LEFT JOIN Brands b          ON f.brand_id    = b.brand_id
WHERE ap.contains_gluten = 0
  AND ap.contains_dairy  = 0
ORDER BY f.food_name;
\end{verbatim}
\end{sqlbox}

\begin{metabox}[Query Translation]
\begin{itemize}
  \item \textbf{English Meaning:} Retrieves all foods that are free from both gluten and dairy,
    together with their category label, nutrition facts, health score, and brand. The allergen
    filter uses an INNER JOIN on ALLERGEN\_PROFILE to guarantee the food has a recorded allergen
    status, while LEFT JOINs on HEALTH\_SCORE, Nutrition\_Metrics, and Brands ensure that foods
    lacking those profiles still appear --- the user is filtering by dietary restriction, not by
    data completeness.
  \item \textbf{Project Feature:} Powers the \textbf{Dietary Filtering} sub-panel in the
    Nutrition Analytics section. Users toggle ``Gluten-Free'' and ``Dairy-Free'' checkboxes;
    the frontend calls \texttt{GET /queries/dietary} and displays matching foods.
  \item \textbf{Criteria Met:}
    \begin{enumerate}[label=\textbf{\arabic*}]
      \item JOIN $\geq$3 tables --- joins \textbf{Foods}, \textbf{FOOD\_CATEGORY},
        \textbf{ALLERGEN\_PROFILE}, \textbf{HEALTH\_SCORE}, \textbf{Nutrition\_Metrics},
        and \textbf{Brands} (6 tables).
      \item OUTER JOIN --- three LEFT JOINs preserve rows that lack health, nutrition, or
        brand data.
    \end{enumerate}
\end{itemize}
\end{metabox}

% ── Query 5 ────────────────────────────────────────────────────────────────────
\subsection{Query 5: Above-Category-Average Calorie Detection}

\begin{sqlbox}[SQL Query]
\small
\begin{verbatim}
SELECT
    f.fdc_id,
    f.food_name,
    fc.category_name AS food_category,
    n.calories,
    n.protein_g,
    n.fat_g,
    n.carbs_g,
    hs.health_score,
    hs.nutriscore_grade,
    b.brand_name
FROM Foods f
JOIN FOOD_CATEGORY fc        ON f.category_id = fc.category_id
JOIN Nutrition_Metrics n     ON f.fdc_id      = n.fdc_id
LEFT JOIN HEALTH_SCORE hs    ON f.fdc_id      = hs.fdc_id
LEFT JOIN Brands b           ON f.brand_id    = b.brand_id
WHERE n.calories > (
    SELECT AVG(n2.calories)
    FROM Foods f2
    JOIN Nutrition_Metrics n2 ON f2.fdc_id = n2.fdc_id
    WHERE f2.category_id = f.category_id
      AND n2.calories IS NOT NULL
)
ORDER BY fc.category_name, n.calories DESC;
\end{verbatim}
\end{sqlbox}

\begin{metabox}[Query Translation]
\begin{itemize}
  \item \textbf{English Meaning:} Identifies food items whose calorie count exceeds the
    average calorie count of all foods in the \emph{same} category. For each food, a
    correlated subquery dynamically computes its category's average calories using
    the \texttt{category\_id} foreign key for efficient correlation, and the outer
    query returns only those above that threshold. This highlights surprisingly energy-dense
    products within each food group.
  \item \textbf{Project Feature:} Drives an \textbf{``Above-Average Calories''} insight
    panel in the Nutrition Analytics section. Users see which products are unusually high
    in calories relative to their category peers, aiding diet planning and nutritional
    awareness.
  \item \textbf{Criteria Met:}
    \begin{enumerate}[label=\textbf{\arabic*}]
      \item JOIN $\geq$3 tables --- outer query joins \textbf{Foods}, \textbf{FOOD\_CATEGORY},
        \textbf{Nutrition\_Metrics}, \textbf{HEALTH\_SCORE}, and \textbf{Brands} (5 tables);
        the subquery joins \textbf{Foods} and \textbf{Nutrition\_Metrics}.
      \item OUTER JOIN --- two LEFT JOINs in the outer query.
      \item Nested SUBQUERY --- the \texttt{WHERE} clause contains a correlated subquery
        that computes the per-category average by referencing the outer \texttt{f.category\_id}.
    \end{enumerate}
\end{itemize}
\end{metabox}

% ── Query 6 ────────────────────────────────────────────────────────────────────
\subsection{Query 6: Nutritional Gap Identification}

\begin{sqlbox}[SQL Query]
\small
\begin{verbatim}
SELECT
    f.fdc_id,
    f.food_name,
    fc.category_name AS food_category,
    dt.type_name     AS data_type,
    b.brand_name,
    n.calories,
    n.protein_g,
    n.fat_g,
    n.carbs_g,
    n.sodium_mg,
    hs.health_score,
    hs.nutriscore_grade
FROM Foods f
JOIN Nutrition_Metrics n     ON f.fdc_id      = n.fdc_id
LEFT JOIN FOOD_CATEGORY fc   ON f.category_id = fc.category_id
LEFT JOIN DATA_TYPE dt       ON f.type_id     = dt.type_id
LEFT JOIN HEALTH_SCORE hs    ON f.fdc_id      = hs.fdc_id
LEFT JOIN Brands b           ON f.brand_id    = b.brand_id
WHERE n.calories   IS NULL
   OR n.protein_g IS NULL
   OR n.fat_g     IS NULL
   OR n.carbs_g   IS NULL
   OR n.sodium_mg IS NULL
ORDER BY f.food_name;
\end{verbatim}
\end{sqlbox}

\begin{metabox}[Query Translation]
\begin{itemize}
  \item \textbf{English Meaning:} Surfaces foods whose Nutrition\_Metrics record contains
    at least one NULL value among the five core nutrient fields (calories, protein, fat,
    carbs, sodium). It identifies data-quality gaps that need attention during the import
    or enrichment pipeline. Category, data type, health, and brand data are brought in via
    LEFT JOINs to provide full context without excluding any incomplete records.
  \item \textbf{Project Feature:} Serves the \textbf{Gap Identification} panel in the
    Nutrition Analytics section. The user clicks ``Identify Gaps,'' the frontend calls
    \texttt{GET /queries/gaps}, and the application lists all foods with incomplete
    nutritional data so the team can target them for enrichment.
  \item \textbf{Criteria Met:}
    \begin{enumerate}[label=\textbf{\arabic*}]
      \item JOIN $\geq$3 tables --- joins \textbf{Foods}, \textbf{Nutrition\_Metrics},
        \textbf{FOOD\_CATEGORY}, \textbf{DATA\_TYPE}, \textbf{HEALTH\_SCORE}, and
        \textbf{Brands} (6 tables).
      \item OUTER JOIN --- four LEFT JOINs preserve context from lookup and profile
        tables even when those rows are missing.
    \end{enumerate}
\end{itemize}
\end{metabox}

% ══════════════════════════════════════════════════════════════════════════════
\section{Criteria Coverage Summary}
% ══════════════════════════════════════════════════════════════════════════════

\begin{table}[H]\centering
\resizebox{\textwidth}{!}{%
\begin{tabular}{lccccc}
\toprule
\textbf{Query} & \textbf{Table Count} & \textbf{(1) JOIN $\geq$3} & \textbf{(2) OUTER JOIN} & \textbf{(3) GROUP BY} & \textbf{(4) SUBQUERY} \\
\midrule
Q1: Full Food Profile      & 7 & \checkmark & \checkmark & & \\
Q2: Multi-Constraint Range & 5 & \checkmark & \checkmark & & \\
Q3: Category Aggregation   & 3 & \checkmark &            & \checkmark & \\
Q4: Dietary Filtering      & 6 & \checkmark & \checkmark & & \\
Q5: Above-Average Calories & 5 & \checkmark & \checkmark & & \checkmark \\
Q6: Gap Identification     & 6 & \checkmark & \checkmark & & \\
\bottomrule
\end{tabular}%
}
\caption{All six queries join at least 3 tables. Each of the four complexity criteria
is covered by at least two queries. Criterion~4 is demonstrated by Q5's correlated subquery
that references the outer \texttt{f.category\_id}. Queries Q1--Q4 and Q6 are currently
implemented in the NutriQuery backend (\texttt{crud.py}); Q5 is planned for the next
release.}
\label{tab:coverage}
\end{table}

\end{document}
